\section{Generate All Binary Strings with No Two Consecutive 1s}
\label{sec:rec-generate-binary-strings}

\problemstatement{Given an integer $n$, generate all binary strings of
length $n$ such that no two $\texttt{1}$s are adjacent (no two
consecutive positions both hold $\texttt{1}$).}

\begin{patternbox}
Another instance of Section~\ref{sec:rec-generate-parentheses}'s
incremental-validity-guard pattern: instead of generating all $2^n$
binary strings and filtering, build the constraint directly into the
choice of what character can be placed \textbf{next}, based only on
\textbf{the single most recently placed character} -- a guard that
depends on a small, fixed amount of local history rather than the entire
string built so far.
\end{patternbox}

\subsection*{Understanding the problem carefully}

At every position, placing a $\texttt{0}$ is always safe, regardless of
what came before -- it can never create two consecutive $\texttt{1}$s.
Placing a $\texttt{1}$ is safe \emph{only if} the immediately preceding
character (if any) was $\texttt{0}$ -- placing it right after another
$\texttt{1}$ would violate the constraint. So the only piece of history
the recursion needs to remember is a single boolean: ``was the last
character placed a $\texttt{1}$?''

\begin{ideabox}[Guard the Choice of $\texttt{1}$ by the Last Character Placed]
$\textsc{Solve}(\textit{current}, \textit{lastWasOne})$: if
$|\textit{current}|=n$: record $\textit{current}$; return. Always:
append $\texttt{0}$; recurse with $\textit{lastWasOne}=\text{false}$;
remove (undo). If $\textit{lastWasOne}$ is \text{false} (i.e.\ the
previous character was not $\texttt{1}$, or this is the first
character): append $\texttt{1}$; recurse with
$\textit{lastWasOne}=\text{true}$; remove (undo).
\end{ideabox}

\subsection*{Why remembering only the single last character is sufficient}

The constraint ``no two consecutive $\texttt{1}$s'' is a purely
\textbf{local} property -- it only ever compares two \emph{adjacent}
positions, never anything further apart. So the only information needed
to decide whether placing a $\texttt{1}$ right now is safe is whether the
immediately preceding position holds a $\texttt{1}$; nothing about
positions further back can affect this decision. This is what allows the
recursion to carry just one bit of state (rather than re-scanning the
entire $\textit{current}$ string at every step to check for violations),
keeping each step's validity check $O(1)$.

\subsection*{Algorithm}

\begin{enumerate}
  \item $\textsc{Solve}(\textit{current}, \textit{lastWasOne})$:
  \begin{enumerate}
    \item If $|\textit{current}|=n$: record $\textit{current}$; return.
    \item Append $\texttt{0}$; recurse with
          $\textit{lastWasOne}=\text{false}$; remove (undo).
    \item If not $\textit{lastWasOne}$: append $\texttt{1}$; recurse
          with $\textit{lastWasOne}=\text{true}$; remove (undo).
  \end{enumerate}
  \item Call $\textsc{Solve}(\texttt{""}, \text{false})$.
\end{enumerate}

\subsection*{Dry run}

\dryrun{Take $n=3$.}

\begin{itemize}
  \item $\textsc{Solve}(\texttt{""}, \text{false})$: place
        $\texttt{0}$: $\textsc{Solve}(\texttt{"0"}, \text{false})$.
  \begin{itemize}
    \item Place $\texttt{0}$: $\textsc{Solve}(\texttt{"00"},
          \text{false})$: place $\texttt{0}$: record \texttt{"000"};
          place $\texttt{1}$ (allowed, last was $0$): record
          \texttt{"001"}.
    \item Place $\texttt{1}$ (allowed): $\textsc{Solve}(\texttt{"01"},
          \text{true})$: place $\texttt{0}$: record \texttt{"010"};
          place $\texttt{1}$ -- \textbf{not allowed} (last was
          $\texttt{1}$), skipped entirely.
  \end{itemize}
  \item Place $\texttt{1}$ (allowed at the very first position):
        $\textsc{Solve}(\texttt{"1"}, \text{true})$: place $\texttt{0}$:
        $\textsc{Solve}(\texttt{"10"}, \text{false})$: place
        $\texttt{0}$: record \texttt{"100"}; place $\texttt{1}$
        (allowed): record \texttt{"101"}. Place $\texttt{1}$ -- not
        allowed (last was $\texttt{1}$), skipped.
\end{itemize}

Recorded strings: \texttt{"000", "001", "010", "100", "101"} -- exactly
the $5$ length-$3$ binary strings with no two consecutive $\texttt{1}$s
(matching the Fibonacci-like count $F(n+2)$ of such strings).

\subsection*{C++ implementation}

\begin{lstlisting}
#include <bits/stdc++.h>
using namespace std;

void solve(string& current, bool lastWasOne, int n, vector<string>& result) {
    if ((int)current.size() == n) {
        result.push_back(current);
        return;
    }
    current.push_back('0');
    solve(current, false, n, result);
    current.pop_back();

    if (!lastWasOne) {
        current.push_back('1');
        solve(current, true, n, result);
        current.pop_back();
    }
}

vector<string> generateBinaryStrings(int n) {
    vector<string> result;
    string current;
    solve(current, false, n, result);
    return result;
}

int main() {
    for (string& s : generateBinaryStrings(3)) cout << s << " ";
    cout << endl;
    return 0;
}
\end{lstlisting}

\begin{complexitybox}
\textbf{Time Complexity.} The number of valid strings of length $n$ is
$O(\phi^n)$ (a Fibonacci-like growth rate, where $\phi$ is the golden
ratio), each costing $O(n)$ to build, giving an overall time complexity
of
\[
  O(n \cdot \phi^n),
\]
substantially smaller than the $O(n \cdot 2^n)$ of generating and
filtering all binary strings.
\textbf{Space Complexity.} $O(n)$ for the recursion depth.
\end{complexitybox}

\begin{takeawaybox}
When a backtracking constraint only ever depends on a small, fixed
amount of recent history (here, just the single previous character),
thread exactly that small amount of state through the recursion's
parameters, rather than re-deriving it by re-scanning the entire
candidate built so far. This keeps each step's validity check $O(1)$ and
generalizes to any constraint expressible in terms of a bounded
``sliding window'' of recent choices.
\end{takeawaybox}
